\documentclass[aspectratio=169,11pt]{ctexbeamer}

\usetheme{Madrid}
\usecolortheme{default}

\usepackage{amsmath,amssymb,mathtools}
\usepackage{bm}
\usepackage{graphicx}

% New lecture folders under "讲义" can keep this style search path.
\makeatletter
\def\input@path{{../模板/}}
\makeatother
\usepackage{yingxing-beamer}

\title{讲义标题}
\subtitle{本讲副标题}
\author{王晨旭}
\date{}

\begin{document}

\begin{frame}
  \titlepage
\end{frame}

\begin{frame}{本讲目标}
\begin{itemize}
  \item 目标一：写清本讲要解决的核心问题。
  \item 目标二：明确学生需要掌握的关键方法。
  \item 目标三：用例题或练习完成迁移训练。
\end{itemize}
\end{frame}

\begin{frame}{核心概念}
\begin{block}{定义或结论}
在这里写本页最重要的定义、定理或方法。
\[
  a_n \to A \quad (n\to\infty)
\]
\end{block}
\end{frame}

\begin{frame}{例题 1}
\textbf{题目：}在这里写例题题干。

\vspace{0.8em}
\textbf{思路：}
\begin{enumerate}
  \item 先判断本题所属模型。
  \item 再给出关键转化或估计。
  \item 最后完成计算或证明。
\end{enumerate}
\end{frame}

\begin{frame}{课堂小结}
\begin{itemize}
  \item 本讲最核心的结论是什么。
  \item 容易出错的条件或边界是什么。
  \item 下一讲需要衔接的问题是什么。
\end{itemize}
\end{frame}

\end{document}
